% Part: normal-modal-logic
% Chapter: filtrations
% Section: euclidean-filtrations

\documentclass[../../../include/open-logic-section]{subfiles}

\begin{document}

\olfileid{nml}{fil}{euc}

\olsection{ترشيحات النماذج الإقليدية}

أما النماذج الإقليدية، وحدها أو مع التسلسل، فلا تكفي فيها طريقة
\olref[acc]{sec}. وبيان ذلك النموذج المرسوم في~\olref{fig:ser-eucl}،
فهو متسلسل وإقليدي. خذ~$\OLId{greek-variable}{Gamma}=\{\OLId{latin-ordinary}{p},\Box \OLId{latin-ordinary}{p}\}$؛ فإذا رُشّح النموذج
عبر~$\OLId{greek-variable}{Gamma}$ اتحدت الفئتان~$[\OLId{latin-ordinary}{w}_\OLMathNumeral{1}]=[\OLId{latin-ordinary}{w}_\OLMathNumeral{3}]$؛ فإن~$\OLId{latin-ordinary}{w}_\OLMathNumeral{1}$ و$\OLId{latin-ordinary}{w}_\OLMathNumeral{3}$
هما وحدهما، من بين أزواج العوالم المختلفة، يتفقان على~$\OLId{greek-variable}{Gamma}$.
ويجب كذلك أن تنتقل أسهم~$\mModel{M}$ إلى كل ترشيح، على نحو ما في
\olref{fig:ser-eucl2}. لكن النموذج الناتج ليس إقليديًا، ولا يمكن جعله
كذلك بزيادة أسهم مع إبقائه ترشيحًا. فإن الإقليدية توجب سهمين متعاكسين
بين~$[\OLId{latin-ordinary}{w}_\OLMathNumeral{2}]$ و$[\OLId{latin-ordinary}{w}_\OLMathNumeral{4}]$، ثم بين~$[\OLId{latin-ordinary}{w}_\OLMathNumeral{2}]$ و$[\OLId{latin-ordinary}{w}_\OLMathNumeral{5}]$.
وهذا ممتنع، لأن~$\Box \OLId{latin-ordinary}{p}$ صادقة عند~$\OLId{latin-ordinary}{w}_\OLMathNumeral{2}$، و$\OLId{latin-ordinary}{p}$ كاذبة عند~$\OLId{latin-ordinary}{w}_\OLMathNumeral{5}$.

\begin{figure}[htpb]
  \centering
  \begin{tikzpicture}[modal]
    \node[world] (w1)
          [label={left: \mFalse{p}},
            label={below: $\mSat{{}}{\Box \OLId{latin-ordinary}{p}}$}] {$\OLId{latin-ordinary}{w}_\OLMathNumeral{1}$};
    \node[world] (w2)
          [label={right: \mTrue{p}},
            label = {below: $\mSat{{}}{\Box \OLId{latin-ordinary}{p}}$},
            right=of w1] {$\OLId{latin-ordinary}{w}_\OLMathNumeral{2}$};
    \draw[->] (w1) to (w2);
    \draw[reflexive above] (w2) to (w2);
    \node[world] (w3)
          [label={left: \mFalse{p}},
            label={below: $\mSat{{}}{\Box \OLId{latin-ordinary}{p}}$},
            below=of w1] {$\OLId{latin-ordinary}{w}_\OLMathNumeral{3}$};
    \node[world] (w4)
          [label={right:\mTrue{p}},
            label={below: $\mSat/{{}}{\Box \OLId{latin-ordinary}{p}}$},
            right=of w3] {$\OLId{latin-ordinary}{w}_\OLMathNumeral{4}$};
    \draw[->] (w3) to (w4);
    \draw[reflexive above] (w4) to (w4);
    \node[world] (w5)
          [label={right: \mFalse{p}},
            label=below:$\mSat/{{}}{\Box \OLId{latin-ordinary}{p}}$,
            right=of w4] {$\OLId{latin-ordinary}{w}_\OLMathNumeral{5}$};
    \draw[->, bend left] (w4) to (w5);
    \draw[->, bend left] (w5) to (w4);
    \draw[reflexive above] (w5) to (w5);
  \end{tikzpicture}
  \caption{نموذج متسلسل وإقليدي.}
  \ollabel{fig:ser-eucl}
\end{figure}

\begin{figure}[ht]
  \centering
  \begin{tikzpicture}[modal,every text node part/.style={align=center}]
    \node[world] (w1)
          [label={left: \mFalse{p}},
            label={right:$[\OLId{latin-ordinary}{w}_\OLMathNumeral{1}]=[\OLId{latin-ordinary}{w}_\OLMathNumeral{3}]$},
            label={below: $\mSat{{}}{\Box \OLId{latin-ordinary}{p}}$}] {$[\OLId{latin-ordinary}{w}_\OLMathNumeral{1}]$};
    \node[world] (w2)
         [label={right: \mTrue{p}},
           label = {below: $\mSat{{}}{\Box \OLId{latin-ordinary}{p}}$},
           right=of w1, yshift=1.5cm] {$[\OLId{latin-ordinary}{w}_\OLMathNumeral{2}]$};
    \draw[->] (w1) to (w2);
    \draw[reflexive above] (w2) to (w2);
    \node[world] (w4)
         [label={right:\mTrue{p}},
           label={below: $\mSat/{{}}{\Box \OLId{latin-ordinary}{p}}$},
           right=of w1,yshift=-1.5cm] {$[\OLId{latin-ordinary}{w}_\OLMathNumeral{4}]$};
    \draw[->] (w1) to (w4);
    \draw[reflexive above] (w4) to (w4);
    \node[world] (w5)
         [label={right: \mFalse{p}},
             label=below:$\mSat/{{}}{\Box \OLId{latin-ordinary}{p}}$,
             right=of w4] {$[\OLId{latin-ordinary}{w}_\OLMathNumeral{5}]$};
    \draw[->, bend left] (w4) to (w5);
    \draw[->, bend left] (w5) to (w4);
    \draw[reflexive above] (w5) to (w5);
  \end{tikzpicture}
  \caption{ترشيح النموذج في~\olref{fig:ser-eucl}.}
  \ollabel{fig:ser-eucl2}
\end{figure}

فلا يكفي، للحصول على ترشيح إقليدي بهذه الطريقة، إغلاق مجموعة اعتباطية
$\OLId{greek-variable}{Gamma}$ تحت أخذ ال!!{formula}s الجزئية. وإنما ننظر في مجموعات
$\OLId{greek-variable}{Gamma}$ \emph{مغلقة جهويًا}، بالمعنى المبين في
\olref[pre]{defn:modallyclosed}. غير أن هذه المجموعات، متى كانت غير
خالية، لا نهائية؛ فلا يُستفاد منها مباشرةً إثبات خاصية النموذج المنتهي،
ولا قابلية النظام المقابل للقرار.

\begin{thm}\ollabel{thm:modal-closed-filt}
  لتكن~$\OLId{greek-variable}{Gamma}$ مغلقة جهويًا، وليكن~$\mModel{M}=\tuple{\OLId{latin-looped}{W},\OLId{latin-looped}{R},\OLId{latin-looped}{V}}$، وليكن
  $\mModel{M^*}=\tuple{\OLId{latin-looped}{W}^*,\OLId{latin-looped}{R}^*,\OLId{latin-looped}{V}^*}$ الترشيح الأخشن لـ~$\mModel{M}$
  عبر~$\OLId{greek-variable}{Gamma}$.
  \begin{enumerate}
  \item إذا كان~$\mModel{M}$ متعديًا، كان~$\mModel{M^*}$ متعديًا.
  \item إذا كان~$\mModel{M}$ إقليديًا، كان~$\mModel{M^*}$ إقليديًا.
  \item إذا كان~$\mModel{M}$ متناظرًا، كان~$\mModel{M^*}$ متناظرًا.
  \end{enumerate}
\end{thm}

\begin{proof}
\begin{enumerate}
  \item لما كان~$\mModel{M^*}$ الترشيح الأخشن، كانت~$\OLId{latin-looped}{R}^*[\OLId{latin-ordinary}{u}][\OLId{latin-ordinary}{v}]$
    مكافئة لـ~$\OLId{latin-looped}{C}_\OLMathNumeral{1}(\OLId{latin-ordinary}{u},\OLId{latin-ordinary}{v})$ بالتعريف. فالتعدي يثبت بأن نفرض~$\OLId{latin-looped}{C}_\OLMathNumeral{1}(\OLId{latin-ordinary}{u},\OLId{latin-ordinary}{v})$
    و$\OLId{latin-looped}{C}_\OLMathNumeral{1}(\OLId{latin-ordinary}{v},\OLId{latin-ordinary}{w})$، ثم نستنتج~$\OLId{latin-looped}{C}_\OLMathNumeral{1}(\OLId{latin-ordinary}{u},\OLId{latin-ordinary}{w})$.
    \iftag{prvBox}{لتكن~$\Box !A \in \OLId{greek-variable}{Gamma}$، وافترض
      $\mSat{\OLId{latin-looped}{M}}{\Box !A}[\OLId{latin-ordinary}{u}]$؛ عندئذ
      $\mSat{\OLId{latin-looped}{M}}{\Box\Box!A}[\OLId{latin-ordinary}{u}]$، لأن~\Ax{4} صالحة في جميع النماذج
      المتعدية؛ ولأن~$\Box\Box!A \in \OLId{greek-variable}{Gamma}$ بالإغلاق، ينتج من
      $\OLId{latin-looped}{C}_\OLMathNumeral{1}(\OLId{latin-ordinary}{u},\OLId{latin-ordinary}{v})$ أن~$\mSat{\OLId{latin-looped}{M}}{\Box!A}[\OLId{latin-ordinary}{v}]$، ثم ينتج من~$\OLId{latin-looped}{C}_\OLMathNumeral{1}(\OLId{latin-ordinary}{v},\OLId{latin-ordinary}{w})$ أن
      $\mSat{\OLId{latin-looped}{M}}{!A}[\OLId{latin-ordinary}{w}]$. }{}\iftag{prvDiamond}{لتكن
      $\Diamond !A \in \OLId{greek-variable}{Gamma}$، وافترض $\mSat{\OLId{latin-looped}{M}}{!A}[\OLId{latin-ordinary}{w}]$؛
      عندئذ~$\mSat{\OLId{latin-looped}{M}}{\Diamond !A}[\OLId{latin-ordinary}{v}]$ بحسب
      $\OLId{latin-looped}{C}_\OLMathNumeral{1}(\OLId{latin-ordinary}{v},\OLId{latin-ordinary}{w})$، لأن~$\Diamond !A \in \OLId{greek-variable}{Gamma}$ بالإغلاق الجهوي. ومن
      $\OLId{latin-looped}{C}_\OLMathNumeral{1}(\OLId{latin-ordinary}{u},\OLId{latin-ordinary}{v})$ نحصل على~$\mSat{\OLId{latin-looped}{M}}{\Diamond\Diamond !A}[\OLId{latin-ordinary}{u}]$، لأن
      $\Diamond\Diamond!A \in \OLId{greek-variable}{Gamma}$ بالإغلاق الجهوي. ولأن
      \Ax{4_\Diamond} صالحة في جميع النماذج المتعدية، يكون
      $\mSat{\OLId{latin-looped}{M}}{\Diamond!A}[\OLId{latin-ordinary}{u}]$.}{}
    % تصحيح برهاني (0459-I1): صُرّح بفرضي العضوية اللازمين لإثبات الشرط C_1.
  \item تمرين. استعمل كون كل من~\Ax{5} و$\Ax{5_\Diamond}$ صالحة في
    جميع النماذج الإقليدية.
  \item تمرين. استعمل كون~\Ax{B} و$\Ax{B_\Diamond}$ صالحة في جميع
    النماذج المتناظرة.
\end{enumerate}
\end{proof}

\begin{prob}
  أكمل برهان~\olref[nml][fil][euc]{thm:modal-closed-filt}.
\end{prob}

\end{document}
